\documentclass[11pt]{amsart}

\usepackage[T1]{fontenc}
\usepackage{lmodern}
\usepackage{microtype}
\usepackage{mathtools,amssymb,amsthm}
\usepackage[hidelinks]{hyperref}

\hypersetup{
  pdftitle={A Ten-Vertex Counterexample to an Endpoint Dichotomy for
Facet Ideals of Graphic Matroids, Answering a Question of DiPasquale,
Fouli, and Kumar}
}

\newtheorem{theorem}{Theorem}
\theoremstyle{remark}
\newtheorem*{remark}{Remark}

\newcommand{\ahat}{\widehat{\alpha}}
\newcommand{\rhat}{\widehat{\rho}}
\newcommand{\rk}{\operatorname{rk}}

\title[A ten-vertex counterexample]
{A Ten-Vertex Counterexample to an Endpoint Dichotomy for Facet Ideals
of Graphic Matroids, Answering a Question of DiPasquale, Fouli, and
Kumar}

\date{}

\subjclass[2020]{13A15, 13F55, 05E40, 05B35}
\keywords{asymptotic resurgence, Waldschmidt constant, facet ideal,
graphic matroid}

\begin{document}

\begin{abstract}
DiPasquale, Fouli, and Kumar asked whether the asymptotic resurgence
of the facet ideal of the graphic matroid of every simple
$2$-connected graph is always one of two extremal values. We give a
ten-vertex graph $W$ for which
\[
c(W)=8,\qquad
\ahat(I(M(W)))=\frac{11}{2},\qquad
\rhat(I(M(W)))=\frac{16}{9}.
\]
Thus
$2-2/c(W)=7/4<16/9<9/5=2-2/|V(W)|$, answering their
Question~7.2 in the negative.
\end{abstract}

\maketitle

For a connected graph $G$, let $M(G)$ be its graphic matroid and let
$I(M(G))\subseteq \Bbbk[x_e:e\in E(G)]$ be its facet ideal. Write
$c(G)$ for the circumference of $G$. DiPasquale, Fouli, and Kumar
proved that every simple $2$-connected graph $G$ satisfies
\[
2-\frac{2}{c(G)}
\le \rhat(I(M(G)))
\le 2-\frac{2}{|V(G)|},
\]
the upper bound being \cite[Theorem~5.5]{DFK}, and asked whether the
asymptotic resurgence is always one of the two endpoints
\cite[Question~7.2]{DFK}. Only the upper bound is used below; the lower
bound is quoted here as context for the question.

Let $W$ be the graph on $\{0,1,\ldots,9\}$ with
\begin{equation}\label{eq:edges}
\begin{split}
E(W)=\{&05,15,06,26,17,27,08,38,48,58,\\
       &29,39,49,79\}.
\end{split}
\end{equation}
Here $ij$ abbreviates the edge $\{i,j\}$. Its graph6 code is
\texttt{I?BDBA[MO}. Deleting vertex $3$ gives the graph in
\cite[Example~6.1]{DFK} under
$(v_1,w_1,u_1,p,q,r,v_2,w_2,u_2)=(8,5,0,4,6,1,9,7,2)$; thus $W$
is obtained from that graph by duplicating one of its degree-two
vertices, $p$.

\begin{theorem}\label{thm:main}
The graph $W$ is simple and $2$-connected, and
\[
c(W)=8,\qquad
\ahat(I(M(W)))=\frac{11}{2},\qquad
\rhat(I(M(W)))=\frac{16}{9}.
\]
Consequently, $W$ is a counterexample to the endpoint dichotomy in
\cite[Question~7.2]{DFK}.
\end{theorem}

\begin{proof}
The graph is simple, and
\[
0-5-1-7-2-6-0,
\quad 0-8-5,
\quad 2-9-7,
\quad 8-3-9,
\quad 8-4-9
\]
is an open-ear decomposition, so $W$ is $2$-connected.

Vertices $3$ and $4$ are nonadjacent twins with neighborhood
$\{8,9\}$. Any cycle containing both is therefore the $4$-cycle
$8-3-9-4-8$, so $W$ has no $10$-cycle.

A $9$-cycle would omit exactly one of $3,4$; by symmetry, suppose it
omits $4$. In $W-4$, the degree-two vertices $1,3,6$ force the paths
\[
5-1-7,\qquad 8-3-9,\qquad 0-6-2.
\]
The remaining cycle edges would form a perfect matching on their six
endpoints. The available edges on those endpoints form the disjoint
triangles
\[
W[\{0,5,8\}]\cong K_3,
\qquad
W[\{2,7,9\}]\cong K_3,
\]
which have no perfect matching. Hence no $9$-cycle exists. Since
\[
5-0-8-3-9-2-7-1-5
\]
is an $8$-cycle, $c(W)=8$.

We next compute the Waldschmidt constant. For
$x\in\mathbb R_{\ge0}^{E(W)}$ and $F\subseteq E(W)$, put
$x(F)=\sum_{e\in F}x_e$. The bond description of the symbolic
polyhedron and its linear-programming formula give
\begin{equation}\label{eq:lp}
\ahat(I(M(W)))
=
\min\bigl\{
 x(E(W)):
 x(B)\ge1\text{ for every bond }B\text{ of }W
\bigr\}
\end{equation}
\cite[Corollary~2.2 and Theorem~2.3]{DFK}.

Because $W$ is $2$-connected, every incident edge cut $\delta(v)$ is
a bond. Put
\[
S=\{3,4,8,9\},
\qquad
B_0=\delta(S)=\{08,58,29,79\}.
\]
Both $W[S]$ and $W[V(W)\setminus S]$ are cycles, so $B_0$ is a bond.
Every edge of $W$ occurs exactly twice in the following weighted sum
of bond inequalities. Thus every feasible $x$ in \eqref{eq:lp}
satisfies
\[
\begin{split}
2x(E(W))
={}&\sum_{v\in\{0,1,2,5,6,7\}}x(\delta(v))
   +2x(\delta(3))+2x(\delta(4))+x(B_0)\\
\ge{}&11.
\end{split}
\]
Hence $\ahat(I(M(W)))\ge 11/2$.

For the reverse inequality, let $J$ be the spanning subgraph formed
by the two cycles
\[
8-0-6-2-7-1-5-8
\qquad\text{and}\qquad
8-3-9-4-8,
\]
of lengths $7$ and $4$. The graph $J$ is connected and bridgeless and
has $7+4=11$ edges. Set
$x_e=1/2$ for $e\in E(J)$ and $x_e=0$ otherwise. If
$B=\delta_W(A)$ is a bond of $W$, then
\[
B\cap E(J)=\delta_J(A).
\]
Since $J$ is connected and bridgeless, $|\delta_J(A)|\ge2$; hence
$x(B)\ge1$. This $x$ is feasible in \eqref{eq:lp} and has total
weight $11/2$. Therefore
\begin{equation}\label{eq:alpha}
\ahat(I(M(W)))=\frac{11}{2}.
\end{equation}

By the identification above,
\[
M(W-3)=M(W)\setminus\{38,39\}.
\]
The graph $W-3$ is the graph of \cite[Example~6.1]{DFK}, whose facet
ideal has asymptotic resurgence $16/9$. Minor monotonicity
\cite[Proposition~2.4]{DFK} therefore gives
\begin{equation}\label{eq:lower}
\rhat(I(M(W)))\ge\frac{16}{9}.
\end{equation}
Since $|V(W-3)|=9$, the value $16/9$ equals $2-2/|V(W-3)|$: the graph
of \cite[Example~6.1]{DFK} sits at the upper endpoint of its own bound
and is not itself a counterexample to Question~7.2.

Since $M(W)$ is simple of rank $9$, the contraction formula
\cite[Theorem~2.5]{DFK} yields
\begin{equation}\label{eq:contraction}
\rhat(I(M(W)))
=
\max\left\{
\frac{9}{\ahat(I(M(W)))},
\max_{e\in E(W)}\rhat(I(M(W)/e))
\right\}.
\end{equation}
The first term is $18/11<16/9$ by \eqref{eq:alpha}.

We use the following consequence of the reduction results in
\cite{DFK}: if $N$ is a loopless graphic matroid of rank $r$, then
\begin{equation}\label{eq:graphic-bound}
\rhat(I(N))\le 2-\frac{2}{r+1}.
\end{equation}
Indeed, simplification preserves asymptotic resurgence, and the
asymptotic resurgence of a direct sum is the maximum over its
matroid-connected components \cite[Theorems~3.2 and~3.3]{DFK}. A
rank-one component contributes $1$. A simple connected graphic
component of rank $s\ge2$ is represented by a simple $2$-connected
graph on $s+1$ vertices, so \cite[Theorem~5.5]{DFK} gives the bound
$2-2/(s+1)\le 2-2/(r+1)$.

For each $e\in E(W)$, the matroid $M(W)/e$ is loopless, graphic, and
of rank $8$. Hence \eqref{eq:graphic-bound} gives
\[
\rhat(I(M(W)/e))\le\frac{16}{9}.
\]
Equation \eqref{eq:contraction} now gives the reverse inequality to
\eqref{eq:lower}, and therefore
\[
\rhat(I(M(W)))=\frac{16}{9}.
\]
The claimed strict inequalities follow from $c(W)=8$ and
$|V(W)|=10$.
\end{proof}

\begin{remark}
DiPasquale, Fouli, and Kumar record that they expect a negative answer
to Question~7.2, while their Sage computations indicate that the answer
is yes for every $2$-connected graph on nine vertices or less
\cite[Table~1 and the discussion following Question~7.2]{DFK}. The
content of Theorem~\ref{thm:main} is therefore the explicit ten-vertex
witness, not a reversal of the source authors' expectation. Concretely,
duplicating $p$ raises the order from $9$ to $10$ and leaves the
asymptotic resurgence at $16/9$, while the upper endpoint moves from
$2-2/9=16/9$ to $2-2/10=9/5$; the value that is extremal for $W-3$ is
therefore strictly interior for $W$. Conditional
on their computation, $W$ has minimum possible order; we do not repeat
that enumeration here. Nothing above determines which values other than
$16/9$ occur strictly between the two endpoints, nor whether $W$ is the
only ten-vertex example.
\end{remark}

\begin{remark}
Every numbered result quoted above is numbered as in \cite{DFK}, and
several of them originate elsewhere. The linear-programming formula for
the Waldschmidt constant, \cite[Theorem~2.3]{DFK}, is stated there with
attribution to Cooper, Embree, H\`a, and Hoefel. The contraction
formula, \cite[Theorem~2.5]{DFK}, is stated there with the citation
\cite[Proposition~3.8 and Theorem~3.9]{DFKmat} attached, \cite{DFKmat}
being the same three authors' companion paper on matroids; minor
monotonicity, \cite[Proposition~2.4]{DFK}, is by contrast a result of
\cite{DFK} itself, proved there from \cite[Corollary~3.4 and
Remark~3.7]{DFKmat}. The remaining items used
above, namely \cite[Corollary~2.2]{DFK},
\cite[Theorems~3.2, 3.3 and~5.5]{DFK}, \cite[Example~6.1]{DFK} and
\cite[Table~1]{DFK}, are from \cite{DFK} itself.
\end{remark}

\begin{thebibliography}{1}

\bibitem{DFK}
M.~DiPasquale, L.~Fouli, and A.~Kumar,
\emph{Asymptotic resurgence of facet ideals of graphic matroids},
arXiv:2607.26294v1 [math.AC], 2026.

\bibitem{DFKmat}
M.~DiPasquale, L.~Fouli, and A.~Kumar,
\emph{Asymptotic resurgence of facet and Stanley--Reisner ideals of
matroids},
arXiv:2607.20892 [math.AC], 2026.

\end{thebibliography}

\end{document}